\documentclass{article}
\usepackage[utf8]{inputenc}
\usepackage[T1]{fontenc}
\usepackage{booktabs}
\usepackage{siunitx}

\\title{Pr\'actica de \'optica: medidas del microscopio}
\date{}

\begin{document}
\maketitle

Las incertidumbres indicadas corresponden a la combinaci\'on de las
incertidumbres de tipo A y B usadas en el script. Los resultados se
escriben como valor nominal(incertidumbre).

\begin{table}[ht]
\centering
\caption{Medidas obtenidas}
\begin{tabular}{lcc}
	oprule
Magnitud & Valor & Unidad \\\\
\midrule
Microscopio 1 & 5(0.001) & mm\\
Microscopio 2 & 145(0.001) & mm \\
Micr\'ometro 1 & 31.5(0.0005) & mm\\
Micr\'ometro 2 & 20(0.0005) & mm \\
\bottomrule
\end{tabular}
\end{table}

El aumento del microscopio se obtiene mediante
\[
M = \frac{d_{\mathrm{microscopio}}}{d_{\mathrm{micr\'ometro}}}.
\]

\begin{table}[ht]
\centering
\caption{Aumento del microscopio}
\begin{tabular}{lcc}
	oprule
Magnitud & Valor & Unidad \\\\
\midrule
Aumento 1 & 0.1587(3.2e-05) & --\\
Aumento 2 & 7.25(0.00019) & -- \\
\bottomrule
\end{tabular}
\end{table}

\end{document}
